一个迭代算法跑到第 \(n\) 轮，本轮的修正量已经小到日志里看不出变化。这轮修正是 \(1/n\) 量级——听起来完全可以忽略。那么整个训练过程累积下来的漂移量呢？也可以忽略吗？

答案是不行。调和级数 \(\sum 1/n\) 的增长虽然极其缓慢，但部分和是无界的——只要步数足够多，积少成多依然可以走向无穷远。反过来，若每轮修正是 \(1/n^2\) 量级，累积漂移就会被严格压在 \(\pi^2/6\) 以内，永远无法越界。两种衰减在曲线变化上几乎难以区分，最终的命运却完全相反。

这正是本章要处理的核心问题：\textbf{一个无限累加的过程，究竟在什么条件下才能产生有限的结果？}

答案不取决于单次增量有多小，而取决于尾部衰减得够不够快。全章的判别工具——比较、积分、比值、根值——回答的都是同一个问题，只是各自拿了不同的尺子去衡量同一段尾巴。

\subsection*{直觉几何化：从``矩形和''到``曲边梯形面积''}

理解离散级数收敛性最直观的视角是\textbf{面积}。将离散求和 \(\sum_{n=1}^N f(n)\) 视为一系列宽度为 \(1\)、高度为 \(f(n)\) 的矩形面积总和；而连续函数积分 \(\int_1^\infty f(x)\,\mathrm{d}x\) 则是曲线下方的连续面积。

只要 \(f(x)\) 单调递减，连续曲线的面积就构成了离散矩形和的天然上下界：
\begin{equation*}
\int_1^{N+1} f(x)\,\mathrm{d}x \;\le\; \sum_{n=1}^N f(n) \;\le\; f(1) + \int_1^N f(x)\,\mathrm{d}x
\end{equation*}

因此，\textbf{级数是否收敛，完全等价于曲线下方的尾部面积是否有限}。

\begin{figure}[htbp]
\centering
\begin{tikzpicture}
\begin{groupplot}[
    group style={group size=2 by 1, horizontal sep=1.5cm},
    width=7.2cm, height=5.5cm,
    xmin=1, xmax=5.5,
    ymin=0, ymax=1.2,
    axis lines=middle,
    xlabel={\(x\)}, ylabel={\(y\)},
    every axis x label/.style={at={(axis description cs:1,0)},anchor=west},
    every axis y label/.style={at={(axis description cs:0,1)},anchor=south},
    xtick={1,2,3,4,5},
    ytick={0.5,1},
    domain=1:5.2,
    samples=100
]

% 左图：1/x (发散)
\nextgroupplot[title={\small (a) \(f(x) = 1/x\)（面积 \(\int_1^\infty \frac{1}{x}\mathrm{d}x = \infty\)）}]
  \fill[red!15] (1,0) rectangle (2,1);
  \fill[red!15] (2,0) rectangle (3,1/2);
  \fill[red!15] (3,0) rectangle (4,1/3);
  \fill[red!15] (4,0) rectangle (5,1/4);
  
  \draw[red!60!black, thick] (1,1) -- (2,1) -- (2,1/2) -- (3,1/2) -- (3,1/3) -- (4,1/3) -- (4,1/4) -- (5,1/4) -- (5,0);
  \addplot[blue, thick, domain=1:5.2] {1/x};
  \node[blue, font=\footnotesize] at (axis cs: 3.2, 0.55) {\(y = \frac{1}{x}\)};
  \node[red!80!black, font=\small] at (axis cs: 3.5, 0.85) {\textbf{面积膨胀 \(\to\) 发散}};

% 右图：1/x^2 (收敛)
\nextgroupplot[title={\small (b) \(f(x) = 1/x^2\)（面积 \(\int_1^\infty \frac{1}{x^2}\mathrm{d}x = 1\)）}]
  \fill[blue!15] (1,0) rectangle (2,1/4);
  \fill[blue!15] (2,0) rectangle (3,1/9);
  \fill[blue!15] (3,0) rectangle (4,1/16);
  \fill[blue!15] (4,0) rectangle (5,1/25);
  
  \draw[blue!60!black, thick] (1,0) -- (1,1/4) -- (2,1/4) -- (2,1/9) -- (3,1/9) -- (3,1/16) -- (4,1/16) -- (4,1/25) -- (5,1/25) -- (5,0);
  \addplot[blue, thick, domain=1:5.2] {1/(x^2)};
  \node[blue, font=\footnotesize] at (axis cs: 2.5, 0.45) {\(y = \frac{1}{x^2}\)};
  \node[blue!80!black, font=\small] at (axis cs: 3.5, 0.85) {\textbf{面积封顶 \(\to\) 收敛}};

\end{groupplot}
\end{tikzpicture}
\caption*{\small 面积视角下的收敛性：\(1/x\) 下方的阴影面积随 \(x \to \infty\) 无限积聚，导致级数发散；而 \(1/x^2\) 下方的总面积被严格限制在有限值以内，将级数部分和封顶。}
\end{figure}

借助积分的换元工具，我们甚至能看清收敛与发散的分界线有多么细微：
\begin{align*}
\int_2^\infty \frac{1}{x \ln x} \,\mathrm{d}x &= \lim_{M\to\infty} \ln(\ln M) = \infty \quad \Rightarrow \quad \text{面积无界（极慢发散）} \\
\int_2^\infty \frac{1}{x \ln^2 x} \,\mathrm{d}x &= \lim_{M\to\infty} \left( \frac{1}{\ln 2} - \frac{1}{\ln M} \right) = \frac{1}{\ln 2} \quad \Rightarrow \quad \text{面积有限（收敛）}
\end{align*}

仅仅多了一个对数因子的二次方，就把无穷远处的尾部面积从无界压回到了有限。

\begin{figure}[H]
\centering
\begin{tikzpicture}
\begin{axis}[
    width=12cm, height=6.5cm,
    xmin=2, xmax=15,
    ymin=0, ymax=0.55,
    axis lines=middle,
    xlabel={\(x\)}, ylabel={\(y\)},
    every axis x label/.style={at={(axis description cs:1,0)},anchor=west},
    every axis y label/.style={at={(axis description cs:0,1)},anchor=south},
    domain=2:15,
    samples=200,
    legend style={at={(0.9,0.88)},anchor=north east,draw=none,fill=white!80}
]

\addplot[red, thick] {1/x};
\addlegendentry{\(y = \frac{1}{x}\)}

\addplot[orange, thick, dashed] {1/(x * ln(x))};
\addlegendentry{\(y = \frac{1}{x\ln x}\)（发散临界）}

\addplot[teal, thick] {1/(x * (ln(x))^2)};
\addlegendentry{\(y = \frac{1}{x\ln^2 x}\)（收敛临界）}

\addplot[blue, thick, dotted] {1/(x^2)};
\addlegendentry{\(y = \frac{1}{x^2}\)}

\node[orange!80!black, font=\small] at (axis cs: 10, 0.16) {发散区域（原函数为 \(\ln(\ln x)\)）};
\node[teal!80!black, font=\small] at (axis cs: 10, 0.03) {收敛区域（原函数为 \(-1/\ln x\)）};

\end{axis}
\end{tikzpicture}
\caption*{\small 临界衰减对比：\(1/(x\ln x)\) 与 \(1/(x\ln^2 x)\) 的图像极难区分，但后者指数上的微小改变，恰好压塌了无穷远处的累积面积。}
\end{figure}

\subsection*{本章路线图}

在接下来的学习中，我们将沿着这条``用不同尺子量尾巴''的主线展开：

\begin{itemize}
  \item \hyperref[sec:02-Calculus/05-Sequences-and-Series/01]{``数列怎样收敛''}：先把``稳定到某个数''写成数列极限，再用单调有界性说明什么时候不必显式求出极限也能确认收敛。
    
  \item \hyperref[sec:02-Calculus/05-Sequences-and-Series/02]{``无穷级数与部分和''}：把``无限次加法''严格定义为部分和序列的极限，并看清通项趋于零为什么仍不足以保证总和有限。
    
  \item \hyperref[sec:02-Calculus/05-Sequences-and-Series/03]{``正项级数的比较''}：正项部分和单调递增，比较、积分、比值与根值判别法分别从不同尺度衡量尾部衰减，也解释学习率条件 \(\sum \alpha_t=\infty\) 与 \(\sum \alpha_t^2<\infty\) 的作用。
    
  \item \hyperref[sec:02-Calculus/05-Sequences-and-Series/04]{``交错与绝对收敛''}：引入符号变化，区分条件收敛与绝对收敛，并说明为什么只有绝对收敛能保证重排不改变和。

  \item \hyperref[sec:02-Calculus/05-Sequences-and-Series/05]{``幂级数是一族函数''}：把常数项推广到变量，确定幂级数的收敛半径，并标出等式成立的有效区域。

  \item \hyperref[sec:02-Calculus/05-Sequences-and-Series/06]{``Taylor 展开与余项：匹配局部变化信息''}：用有限阶展开逼近函数，再用余项界说明局部模型在多大范围内仍值得信任。
\end{itemize}

读完本章，当你再次面对一个无限求和时，直觉将驱动你首先去考察尾部的衰减尺度；而看到一个近似展开时，你的注意力也将自然落到余项的控制范围上。